\documentclass[11pt]{article}

\usepackage[T1]{fontenc}
\usepackage[utf8]{inputenc}
\usepackage[a4paper,margin=1in]{geometry}
\usepackage{graphicx}
\usepackage{booktabs}
\usepackage{amsmath}
\usepackage{amssymb}
\usepackage{xcolor}
\usepackage{titlesec}
\usepackage{caption}
\usepackage[colorlinks=true,linkcolor=navy,urlcolor=navy,citecolor=navy]{hyperref}

\definecolor{navy}{HTML}{1F3864}
\definecolor{accent}{HTML}{2E7D5B}
\definecolor{grey}{HTML}{595959}

\titleformat{\section}{\normalfont\Large\bfseries\color{navy}}{\thesection}{0.6em}{}[{\color{accent}\titlerule[1pt]}]
\titleformat{\subsection}{\normalfont\large\bfseries\color{navy}}{\thesubsection}{0.5em}{}
\captionsetup{font=small,labelfont={bf,color=navy}}

\begin{document}

% ================= Title =================
\begin{center}
  {\LARGE\bfseries\color{navy} Expected Goals (xG) Modelling in Football}\\[6pt]
  {\large\itshape\color{grey} A Mathematical Modelling Approach}\\[4pt]
  {\small\color{grey} ACM 40960 \ $\cdot$ \ Mathematical Modelling}
\end{center}
\vspace{2pt}
{\color{navy}\hrule height 1.2pt}
\vspace{12pt}

% ================= Abstract =================
\noindent\textbf{Abstract.} This project builds and studies a calibrated Expected
Goals (xG) model for football using the StatsBomb open dataset. The work treats
xG as a mathematical model rather than only a prediction task. The two main shot
features, distance and angle, are derived from pitch geometry, the baseline
logistic model is framed as a maximum likelihood optimisation problem, and the
fitted model is then extended in three directions: a Monte Carlo match simulation
based on the Poisson-binomial distribution, a Bayesian version that quantifies
uncertainty, and a sensitivity analysis. A simple distance and angle model
reaches an AUC of 0.75, and adding contextual features raises this to 0.78.
Random Forest and XGBoost match the enhanced logistic model almost exactly, which
shows that with these features the choice of algorithm matters far less than the
information available. All models are well calibrated. The simulation recovers
known real-world results, identifying Leicester City and Atl\'etico Madrid as the
largest over-performers of the 2015/16 season.

\vspace{6pt}
\tableofcontents
\vspace{10pt}

% ================= Introduction =================
\section{Introduction}

Expected Goals (xG) is a way of measuring how likely a shot is to end up as a
goal. Each shot is given a value between 0 and 1, where a higher value means a
better chance of scoring, based on features such as the distance from goal, the
angle to goal, the body part used, and the amount of defensive pressure. Over the
last decade xG has become a standard metric in football analytics and is now used
widely by clubs, broadcasters, and researchers. Because goals are rare, a measure
that estimates the quality of a shot rather than just counting goals gives a more
stable picture of performance.

The aim of this project is not only to fit an accurate xG model but to treat it
as a mathematical object. The shot geometry is derived from first principles, the
fitting of the baseline model is framed as an optimisation problem, the model is
extended to whole match outcomes through simulation, and the uncertainty in the
predictions is quantified. The work is carried out entirely on open data so that
it remains reproducible.

% ================= Literature Review =================
\section{Literature Review}

Early xG models were kept deliberately simple. Rathke (2017) showed that a
logistic regression model based mainly on shot distance and angle could already
give a reasonable estimate of goal probability, which set the basic template that
most later work has built on. Much of the later research focused on adding more
information and using more flexible methods. Lucey et al.\ (2015) showed that
strategic and spatial features improved prediction beyond shot location alone, and
Anzer and Bauer (2021) combined synchronised positional and event data to build
one of the more detailed models in the field. Mead et al.\ (2023) argued that
standard models leave out useful information such as player and team ability, and
researchers have increasingly moved from logistic regression toward tree based
ensemble methods such as Random Forest and gradient boosting.

A separate strand of work has focused on interpretability and on adjusting models
for who takes the shot. Hewitt and Karaku\c{s} (2023) trained position adjusted
models, and Scholtes and Karaku\c{s} (2024) used a Bayesian hierarchical approach
to correct for player and position. A common use of xG is to judge finishing
ability by comparing a player's actual goals to their total xG. This use has been
questioned: Davis and Robberechts (2024) found that a player's overperformance is
not stable from one season to the next, and that xG models can carry biases
because stronger players and teams take more shots. Their work is a useful
reminder that a finishing table should be read with care rather than treated as a
clean measure of skill.

xG is also a natural building block for modelling whole match outcomes. A long
line of work models the number of goals a team scores using Poisson based models,
starting with Dixon and Coles (1997), and later placed in a Bayesian hierarchical
framework by Baio and Blangiardo (2010). Because each shot already has a goal
probability, the goals in a match can be treated as a sum of independent Bernoulli
trials, which links the shot level model to a simulation of results at the match
level. The present project draws these threads together on open data, with a
particular emphasis on calibration and on the mathematical structure behind each
step.

% ================= Data =================
\section{Data}

The data source is the StatsBomb open dataset. In the open release only the
2015/16 season is fully covered for the major leagues, so the four full 2015/16
seasons of La Liga, the Premier League, Serie A, and Ligue 1 were combined into a
single dataset. This gives a large and unbiased sample rather than the
player-specific subsets available for other seasons. After removing penalties,
which are always taken from the same spot and would distort a distance and angle
model, the dataset contains 37{,}488 shots, of which 3{,}569 were goals, an
overall conversion rate of 9.5 percent. Each shot records its pitch coordinates,
body part, technique, shot type, play pattern, whether the shooter was under
pressure, and whether it was a first-time shot.

% ================= Feature engineering: geometry =================
\section{Shot Geometry}

The pitch is 120 units long and 80 units wide. The goal being attacked lies on
the line $x = 120$, with posts at $y = 36$ and $y = 44$, so the goal is $w = 8$
units wide and its centre is $C = (120, 40)$. Two features are derived from the
raw location $(x, y)$.

The distance is the straight-line distance to the centre of the goal,
\begin{equation}
  d = \sqrt{(120-x)^2 + (40-y)^2}.
\end{equation}
The angle is the width of the goal as seen from the shot, that is, the angle
$\theta$ subtended at $P=(x,y)$ by the two posts $A=(120,36)$ and $B=(120,44)$.
With $p = |PA|$ and $q = |PB|$ the distances to each post, the law of cosines
applied to triangle $APB$, whose side opposite $\theta$ is the goal of length
$w=8$, gives
\begin{equation}
  \theta = \arccos\!\left(\frac{p^2 + q^2 - w^2}{2pq}\right).
\end{equation}

Figure~\ref{fig:geom} shows lines of equal distance and equal angle on the pitch.
Each equal-angle curve is a circular arc through both posts, a consequence of the
inscribed angle theorem. Figure~\ref{fig:validation} confirms that both features
carry genuine signal: conversion falls sharply with distance and rises with a
wider angle, which justifies using $d$ and $\theta$ as the inputs to the baseline
model.

\begin{figure}[h]
  \centering
  \includegraphics[width=\textwidth]{geometry_iso_curves.png}
  \caption{Lines of equal distance (left) and equal angle (right). The red segment is the goal.}
  \label{fig:geom}
\end{figure}

\begin{figure}[h]
  \centering
  \includegraphics[width=\textwidth]{feature_validation.png}
  \caption{Empirical conversion rate against distance and angle, confirming both features are predictive.}
  \label{fig:validation}
\end{figure}

% ================= Baseline model =================
\section{Baseline Model: Logistic Regression}

For each shot $i$ with distance $d_i$ and angle $\theta_i$ and outcome
$y_i \in \{0,1\}$, the probability of a goal is modelled as
\begin{equation}
  p_i = \sigma(\eta_i) = \frac{1}{1 + e^{-\eta_i}}, \qquad
  \eta_i = \beta_0 + \beta_1 d_i + \beta_2 \theta_i,
\end{equation}
which is linear on the log-odds scale,
$\log\!\big(p_i/(1-p_i)\big) = \beta_0 + \beta_1 d_i + \beta_2 \theta_i$.

\paragraph{Fitting as an optimisation problem.}
Each outcome is a Bernoulli trial, so the log-likelihood is
\begin{equation}
  \ell(\boldsymbol{\beta}) = \sum_{i=1}^{n}\big[\,y_i \log p_i + (1-y_i)\log(1-p_i)\,\big],
\end{equation}
and the maximum likelihood estimate is
$\hat{\boldsymbol{\beta}} = \arg\max_{\boldsymbol{\beta}} \ell(\boldsymbol{\beta})$.
Setting the gradient to zero gives the score equations
$\mathbf{X}^\top(\mathbf{y} - \mathbf{p}) = \mathbf{0}$, which have no closed-form
solution and are solved numerically by Newton--Raphson. The log-likelihood is
concave, so the solution is the unique global maximum.

The fitted coefficients were $\beta_0 = -1.72$, $\beta_1 = -0.073$ for distance,
and $\beta_2 = 0.023$ for angle, all highly significant. As odds ratios, each
extra unit of distance multiplies the odds of scoring by $0.93$, a reduction of
about 7 percent, and each extra degree of angle multiplies them by $1.02$. The
resulting xG surface is shown in Figure~\ref{fig:surface}.

\begin{figure}[h]
  \centering
  \includegraphics[width=0.62\textwidth]{baseline_xg_surface.png}
  \caption{Baseline xG surface over the attacking third, brightest in front of goal.}
  \label{fig:surface}
\end{figure}

% ================= Evaluation =================
\section{Evaluation}

Models are assessed on a held-out test set of 20 percent of shots, using three
measures. The area under the ROC curve (AUC) is the probability that a random goal
is given a higher xG than a random non-goal, measuring ranking only. The Brier
score, $\mathrm{BS} = \tfrac{1}{n}\sum_i (p_i - y_i)^2$, measures the accuracy of
the probabilities, with lower being better. Calibration is checked with a
reliability plot and summarised by the expected calibration error (ECE), the
average gap between predicted probability and observed rate.

The baseline reached an AUC of 0.746 and a Brier score of 0.0795, beating the
no-skill reference of 0.086. Its reliability curve
(Figure~\ref{fig:baseeval}) sits almost exactly on the diagonal, so its
probabilities can be read as genuine chances.

\begin{figure}[h]
  \centering
  \includegraphics[width=\textwidth]{baseline_evaluation.png}
  \caption{Baseline model: ROC curve (left) and calibration plot (right).}
  \label{fig:baseeval}
\end{figure}

% ================= Contextual features and enhanced model =================
\section{Contextual Features and the Enhanced Model}

The descriptive columns were turned into numeric features by one-hot encoding,
where each category becomes a 0/1 column. For each feature one level was left out
as a reference (right foot, normal technique, open play, regular play), so every
coefficient reads as a comparison to a typical shot. Figure~\ref{fig:cats} shows
how conversion varies across these categories. Headers and shots from a counter
attack convert relatively well, and first-time shots convert better than shots
where the player takes a touch first.

\begin{figure}[h]
  \centering
  \includegraphics[width=\textwidth]{category_conversion.png}
  \caption{Empirical conversion rate by body part, technique, play pattern, and shot context.}
  \label{fig:cats}
\end{figure}

Fitting the logistic model on all 20 features raised the AUC to 0.777 and improved
the Brier score to 0.0772. The coefficients, shown as odds ratios in
Figure~\ref{fig:or}, are revealing. Headers now have the strongest negative
effect, an odds ratio near $0.22$, which appears to contradict their high raw
conversion. The explanation is that headers convert well only because they occur
close to goal, and once distance and angle are already accounted for, a header
from the same position is clearly worse than a foot shot. This is the model
separating location from technique, and it illustrates why a multi-feature model
gives a fairer picture. Shots under pressure score less often, and shots from a
counter attack score more.

\begin{figure}[h]
  \centering
  \includegraphics[width=0.8\textwidth]{enhanced_odds_ratios.png}
  \caption{Effect of each feature in the enhanced model. Bars right of 1 raise the odds of scoring, bars left of 1 lower them.}
  \label{fig:or}
\end{figure}

% ================= Tree models and comparison =================
\section{Tree Models and Comparison}

Two tree based models were fitted with the same features, split, and evaluation.
A Random Forest grows many decision trees on random slices of the data and
averages them, while XGBoost builds trees one after another, each correcting the
errors of the last. Both were tuned lightly using a validation slice of the
training data.

\begin{table}[h]
  \centering
  \caption{Model comparison on the held-out test set. Arrows show the good direction.}
  \label{tab:compare}
  \begin{tabular}{lcccc}
    \toprule
    Model & AUC $\uparrow$ & Brier $\downarrow$ & LogLoss $\downarrow$ & ECE $\downarrow$ \\
    \midrule
    Baseline logit (distance + angle) & 0.746 & 0.0795 & 0.280 & 0.0162 \\
    Enhanced logit (context)          & 0.777 & 0.0772 & 0.270 & 0.0138 \\
    Random Forest                     & 0.779 & 0.0769 & 0.269 & 0.0126 \\
    XGBoost                           & 0.780 & 0.0769 & 0.269 & 0.0114 \\
    \midrule
    StatsBomb xG (benchmark)          & 0.816 & 0.0714 & 0.251 & 0.0116 \\
    \bottomrule
  \end{tabular}
\end{table}

Table~\ref{tab:compare} and Figure~\ref{fig:compare} tell a clear story. The
largest single gain comes from adding contextual features to the baseline. After
that, the enhanced logistic model, the Random Forest, and XGBoost are almost
indistinguishable on every measure. This means the models have reached the ceiling
of what the available features can explain, and the remaining gap to the StatsBomb
benchmark is due to information, chiefly the positions of defenders and the
goalkeeper, rather than to model choice. All models are well calibrated, and on
calibration error the XGBoost model matches the professional benchmark. Applying
Platt scaling or isotonic regression to the already-calibrated XGBoost model did
not improve test-set calibration, so no correction was retained.

\begin{figure}[h]
  \centering
  \includegraphics[width=\textwidth]{model_comparison.png}
  \caption{AUC, Brier score, and calibration across all models.}
  \label{fig:compare}
\end{figure}

% ================= Applied outputs =================
\section{Applied Outputs}

XGBoost was used as the final model for the applied outputs, being the strongest
on every measure and the best calibrated. To avoid a model scoring shots it was
trained on, each shot was assigned an xG using five-fold cross-validation. As a
whole-dataset check, the total xG of 3{,}568 almost exactly matched the 3{,}569
actual goals, confirming good calibration across the entire sample.

Figure~\ref{fig:shotmap} maps shots coloured by xG, with the highest values
concentrated centrally in the box. Finishing was then measured for each player as
actual goals minus total xG. Figure~\ref{fig:finish} shows the biggest over and
under-performers among players with at least 40 shots. The over-performers are led
by elite forwards of the era such as Su\'arez, Ibrahimovic, Bale, Higua\'in, and
Griezmann, which is strong evidence that the model behaves sensibly. As the
literature review noted, a single season of finishing is noisy and biased, so this
table describes who out-performed their chances in this season rather than a
permanent measure of skill.

\begin{figure}[h]
  \centering
  \includegraphics[width=0.82\textwidth]{shot_map_xg.png}
  \caption{A sample of shots coloured by predicted xG.}
  \label{fig:shotmap}
\end{figure}

\begin{figure}[h]
  \centering
  \includegraphics[width=0.78\textwidth]{finishing_table.png}
  \caption{Finishing: goals above or below expectation, players with at least 40 shots.}
  \label{fig:finish}
\end{figure}

% ================= Match simulation =================
\section{Match Simulation}

Once every shot has an xG value $p_i$, a shot can be treated as an independent
Bernoulli trial. A team's goals in a match, $G = \sum_i B_i$ with
$B_i \sim \mathrm{Bernoulli}(p_i)$, then follow a Poisson-binomial distribution
with mean $\sum_i p_i$ and variance $\sum_i p_i(1-p_i)$. Treating the two teams'
goals as independent, the match outcome probabilities are
$P(\text{home win}) = P(G_H > G_A)$ and so on, and the expected points for the
home team are
\begin{equation}
  \mathbb{E}[\text{pts}_H] = 3\,P(G_H > G_A) + 1\,P(G_H = G_A).
\end{equation}
These quantities were estimated by Monte Carlo, simulating every match 10{,}000
times. Figure~\ref{fig:match} validates the method on the highest-shot-volume
match: the simulated goal distribution matches the exact Poisson-binomial
distribution, and the outcome probabilities show that Juventus deserved to win 71
percent of the time yet only drew.

\begin{figure}[h]
  \centering
  \includegraphics[width=\textwidth]{example_match_sim.png}
  \caption{Example match. Left: Monte Carlo goal counts (dots) against the exact Poisson-binomial distribution (bars). Right: outcome probabilities.}
  \label{fig:match}
\end{figure}

Summing expected points across a season gives an xG-based league table.
Figure~\ref{fig:xpoints} compares actual points to expected points for all four
leagues. Teams above the diagonal earned more real points than their chances
deserved. The two largest over-performers were Leicester City ($+20$) and
Atl\'etico Madrid ($+21$), which matches the well-known story of the 2015/16
season, when Leicester won the Premier League against expectation and Atl\'etico
were famed for their efficiency.

\begin{figure}[h]
  \centering
  \includegraphics[width=0.72\textwidth]{xpoints_vs_actual.png}
  \caption{Actual points against expected points from xG. Teams above the line over-performed their chances.}
  \label{fig:xpoints}
\end{figure}

% ================= Bayesian model =================
\section{Bayesian Model and Uncertainty}

The models above give a single estimate for each quantity. A Bayesian treatment
produces a full posterior distribution and therefore credible intervals. Keeping
the same likelihood and placing weakly-informative priors
$\alpha \sim \mathcal{N}(0,5^2)$ and $\beta_k \sim \mathcal{N}(0,5^2)$, the
posterior $p(\alpha,\boldsymbol{\beta}\mid\mathbf{y})$ was sampled. Because the
dataset is large, the Bernstein--von Mises theorem ensures the posterior is close
to $\mathcal{N}(\hat{\boldsymbol{\beta}}, \widehat{\boldsymbol{\Sigma}})$, which
was used for the results here and reproduced with full Hamiltonian Monte Carlo in
Stan.

The posterior for the distance coefficient had a 95 percent credible interval of
about $[-0.085, -0.068]$ and the angle coefficient $[0.018, 0.024]$, both well
away from zero. Figure~\ref{fig:bayes} shows the coefficient posteriors and a
credible band on the xG versus distance curve. The band is wider close to goal,
where fewer shots are taken and the model is less certain, and narrower far out.
This is uncertainty quantification made explicit: any prediction can now be
reported with a credible interval.

\begin{figure}[h]
  \centering
  \includegraphics[width=\textwidth]{bayes_uncertainty.png}
  \caption{Posterior distributions of the coefficients (left) and xG against distance with a 95 percent credible band (right).}
  \label{fig:bayes}
\end{figure}

% ================= Sensitivity =================
\section{Sensitivity Analysis}

The behaviour of the model was studied both locally and globally. Locally, the
derivatives of the logistic model follow from
$\sigma'(\eta) = \sigma(\eta)(1-\sigma(\eta))$:
\begin{equation}
  \frac{\partial p}{\partial d} = \beta_1\, p(1-p),
  \qquad
  \frac{\partial p}{\partial \theta} = \beta_2\, p(1-p).
\end{equation}
Both share the factor $p(1-p)$, which is largest when $p = 0.5$. The model is
therefore most sensitive to position for shots of moderate xG, physically a ring
around the edge of the six-yard box, and least sensitive for near-certain or
near-hopeless chances. This is visible in Figure~\ref{fig:senslocal}.

\begin{figure}[h]
  \centering
  \includegraphics[width=\textwidth]{sensitivity_local.png}
  \caption{Local sensitivity: change in xG per unit distance and per degree of angle (left), and where on the pitch xG is most sensitive to position (right).}
  \label{fig:senslocal}
\end{figure}

Globally, variance-based Sobol indices measure how much of the variation in xG
across all shots is driven by each input. Distance had a first-order index of 0.62
and angle 0.31, so distance is the dominant driver with angle second. The small
gap between the first-order and total-effect indices (0.62 against 0.69 for
distance) indicates only a modest interaction between the two features
(Figure~\ref{fig:sobol}).

\begin{figure}[h]
  \centering
  \includegraphics[width=0.6\textwidth]{sensitivity_sobol.png}
  \caption{Global sensitivity: first-order and total-effect Sobol indices.}
  \label{fig:sobol}
\end{figure}

% ================= Discussion =================
\section{Discussion and Limitations}

The results support three main conclusions. First, most of the model's predictive
power comes from a small set of geometric and contextual features, and the largest
single improvement is the move from distance and angle alone to the full feature
set. Second, once these features are fixed, flexible models offer almost no
advantage over a well-specified logistic regression, so the interpretable model is
a reasonable final choice. Third, the gap to a professional model is explained by
missing information rather than by method, which is a useful and honest boundary on
what open event data can achieve.

Several limitations should be noted. The data lack freeze-frame information on
defender and goalkeeper positions, which is the single largest missing input and
the main reason for the gap to the StatsBomb benchmark. Only one fully covered
season was available, so the model cannot be tested across seasons. The finishing
table is based on a single season and, as the literature shows, such measures are
noisy and biased, so they should not be read as fixed skill. The match simulation
treats the two teams' goals as independent and ignores game state, and it uses
open-play, non-penalty shots only, so simulated scorelines omit penalties and own
goals. Finally, the global sensitivity analysis assumes distance and angle vary
independently and uniformly, whereas in real shots they are correlated.

% ================= Conclusion =================
\section{Conclusion}

This project built a calibrated xG model on open data and treated it as a
mathematical model throughout. The shot geometry was derived from first
principles, the baseline fit was framed as a maximum likelihood optimisation
problem, and the model was extended to match outcomes through a Poisson-binomial
Monte Carlo simulation, to uncertainty through a Bayesian treatment, and to
behaviour through a sensitivity analysis. The final model is accurate, well
calibrated, and interpretable, and its outputs recover known real-world results.
The clearest route to further improvement would be richer positional data, which
the analysis identifies as the main factor separating this model from a
professional one.

% ================= References =================
\section*{References}
\addcontentsline{toc}{section}{References}
\begin{itemize}\setlength{\itemsep}{2pt}
  \item Anzer, G., \& Bauer, P. (2021). A goal scoring probability model for shots based on synchronized positional and event data in football (soccer). \textit{Frontiers in Sports and Active Living}, 3, 624475.
  \item Baio, G., \& Blangiardo, M. (2010). Bayesian hierarchical model for the prediction of football results. \textit{Journal of Applied Statistics}, 37(2), 253--264.
  \item Davis, J., \& Robberechts, P. (2024). Biases in expected goals models confound finishing ability. \textit{arXiv preprint} arXiv:2401.09940.
  \item Dixon, M. J., \& Coles, S. G. (1997). Modelling association football scores and inefficiencies in the football betting market. \textit{Journal of the Royal Statistical Society: Series C (Applied Statistics)}, 46(2), 265--280.
  \item Hewitt, J. H., \& Karaku\c{s}, O. (2023). A machine learning approach for player and position adjusted expected goals in football (soccer). \textit{Franklin Open}, 4, 100034.
  \item Lucey, P., Bialkowski, A., Monfort, M., Carr, P., \& Matthews, I. (2015). Quality vs quantity: Improved shot prediction in soccer using strategic features from spatiotemporal data. \textit{MIT Sloan Sports Analytics Conference}.
  \item Mead, J., O'Hare, A., \& McMenemy, P. (2023). Expected goals in football: Improving model performance and demonstrating value. \textit{PLoS ONE}, 18(4), e0282295.
  \item Rathke, A. (2017). An examination of expected goals and shot quality from a goalkeeper's perspective. \textit{International Journal of Computer Science in Sport}, 16(2), 69--79.
  \item Scholtes, A., \& Karaku\c{s}, O. (2024). Bayes-xG: player and position correction on expected goals using a Bayesian hierarchical approach. \textit{Frontiers in Sports and Active Living}, 6, 1348983.
  \item StatsBomb (2024). StatsBomb open data. \url{https://github.com/statsbomb/open-data}
\end{itemize}

\end{document}
